\section{System Model}
\label{sec:system-model}

\subsection{Received signal}

Consider a hybrid receiver with $N_a$ antennas, $N_{rf}$ radio frequency
chains, and an OFDM waveform of $N$ subcarriers with sampling interval $T_s$.
Write $\vect{s} \in \mathbb{C}^{N}$ for the transmitted frequency-domain
symbols, $\vect{h} \in \mathbb{C}^{L}$ for the $L$-tap combined channel after
analog beamforming, and $\mat{F} \in \mathbb{C}^{N \times L}$ for the first $L$
columns of the unitary discrete Fourier transform matrix. Let
$\phi_n$, $n = 0, \dots, N-1$, denote the oscillator phase during sample $n$ of
the symbol. The time-domain received samples are
\begin{equation}
  \label{eq:time-domain}
  y_n = e^{j \phi_n} \sum_{\ell=0}^{L-1} h_\ell \, x_{n-\ell} + w_n ,
  \qquad w_n \sim \mathcal{CN}(0, \sigma^2) ,
\end{equation}
where $\vect{x} = \mat{W}^{\mathsf{H}} \vect{s}$ is the transmitted time-domain
block. Because the phase multiplies the signal after combining, it is a scalar
sequence rather than a per-antenna one, which is what makes the hybrid case
harder than the digital case and also what makes joint estimation tractable.

Transforming Eq.~\eqref{eq:time-domain} to the frequency domain gives the
familiar decomposition into a common rotation and a leakage term,
\begin{align}
  \label{eq:freq-domain}
  Y_k
  &= \underbrace{J_0 H_k s_k}_{\text{common}}
     + \underbrace{\sum_{m \ne k} J_{k-m} H_m s_m}_{\text{ICI}}
     + W_k , \\
  \label{eq:jdef}
  J_q &= \frac{1}{N} \sum_{n=0}^{N-1} e^{j \phi_n} e^{-j 2\pi q n / N} ,
\end{align}
with $H_k$ the channel frequency response at subcarrier $k$. Equation
\eqref{eq:jdef} defines the phase-noise spectral coefficients, and every
approximation in the literature amounts to a claim about how quickly $J_q$
decays in $q$.

\subsection{Leakage statistics without the small-phase approximation}

The usual step is to write $e^{j\phi_n} \approx 1 + j \phi_n$ and propagate a
Gaussian $\phi_n$ through the linearisation. That approximation fails once the
accumulated phase variance over a symbol approaches one radian, which at
$\beta = 200$~kHz and $N T_s = 3.2~\mu$s it does. We instead use the exact
characteristic function of the Wiener phase process. For a free-running
oscillator, $\phi_n = \phi_{n-1} + u_n$ with
$u_n \sim \mathcal{N}(0, 2\pi \beta T_s)$, so for any lag $\Delta$
\begin{align}
  \label{eq:charfun}
  \mathbb{E}\!\left[ e^{j(\phi_n - \phi_{n-\Delta})} \right]
    &= e^{-\pi \beta T_s |\Delta|} , \\
  \label{eq:ici-cov}
  \mathbb{E}\!\left[ J_q J_{q'}^{*} \right]
    &= \frac{1}{N^2} \sum_{n,n'} e^{-\pi \beta T_s |n - n'|}
       e^{-j 2\pi (qn - q'n')/N} .
\end{align}
Equation~\eqref{eq:ici-cov} is exact, not a second-order expansion, and it
holds for any linewidth. Collecting the coefficients into
$\vect{J} = [J_{-N/2}, \dots, J_{N/2-1}]^{\mathsf{T}}$ and writing
$\mat{C} = \mathbb{E}[\vect{J} \vect{J}^{\mathsf{H}}]$, the sum in
Eq.~\eqref{eq:ici-cov} is a doubly indexed geometric series and evaluates in
closed form to
\begin{equation}
  \label{eq:cov-closed}
  [\mat{C}]_{q,q'} =
  \frac{1}{N^2} \cdot
  \frac{\bigl(1 - e^{-2\pi\beta T_s N}\bigr) \, \gamma^{|q-q'|}}
       {\bigl(1 - \gamma e^{j 2\pi (q-q')/N}\bigr)
        \bigl(1 - \gamma e^{-j 2\pi (q-q')/N}\bigr)} ,
\end{equation}
with $\gamma = e^{-\pi \beta T_s}$. The matrix is Hermitian Toeplitz and its
eigenvalues decay geometrically away from the main diagonal.

\begin{proposition}[Numerical rank of the leakage covariance]
\label{prop:rank}
Fix a tolerance $\epsilon$. The number of eigenvalues of $\mat{C}$ exceeding
$\epsilon \tr(\mat{C})$ satisfies
$R(\epsilon) \le 2 \lceil \beta T_s N \rceil + 1 +
 O\!\bigl(\log(1/\epsilon)\bigr)$.
\end{proposition}

Proposition~\ref{prop:rank} is the structural fact the estimator exploits. At
our operating point, $\beta T_s N = 0.64$ and $R(10^{-3}) = 9$ against
$N = 1024$. The leakage lives in a nine-dimensional subspace, and an estimator
that models that subspace recovers almost all of the energy that a least
squares estimator discards.

\subsection{Signal chain}

Figure~\ref{fig:chain} shows where each quantity enters. The single local
oscillator feeding the mixer is the source of the common phase process, and the
estimator sits after the fast Fourier transform, so it observes
Eq.~\eqref{eq:freq-domain} directly.

\begin{figure*}[t]
  \centering
  \begin{tikzpicture}[
      node distance=6mm and 7mm,
      blk/.style={draw,rounded corners=1.5pt,align=center,font=\scriptsize,
                  minimum height=8.5mm,minimum width=15mm,inner xsep=3pt},
      osc/.style={draw,circle,align=center,font=\scriptsize,minimum size=10mm},
      flow/.style={-{Stealth[length=1.8mm]},semithick},
    ]
    \node[blk,align=center] (arr) {Antenna\\array $N_a$};
    \node[blk,right=of arr] (ps) {Analog\\combiner};
    \node[blk,right=of ps] (lna) {LNA};
    \node[blk,right=of lna] (mix) {Mixer};
    \node[blk,right=of mix] (adc) {ADC};
    \node[blk,right=of adc] (fft) {Remove CP\\and FFT};
    \node[blk,right=of fft,fill=black!6] (est) {\est{}\\estimator};
    \node[blk,right=of est] (eq) {Equalizer\\and demapper};

    \node[osc,below=9mm of mix] (lo) {LO};
    \node[font=\scriptsize,below=1.5mm of lo,align=center]
      {free running\\$\beta = 200$\,kHz};

    \draw[flow] (arr) -- (ps);
    \draw[flow] (ps) -- node[above,font=\scriptsize] {$N_{rf}$} (lna);
    \draw[flow] (lna) -- (mix);
    \draw[flow] (mix) -- (adc);
    \draw[flow] (adc) -- node[above,font=\scriptsize] {$y_n$} (fft);
    \draw[flow] (fft) -- node[above,font=\scriptsize] {$Y_k$} (est);
    \draw[flow] (est) -- node[above,font=\scriptsize] {$\hat{H}_k$} (eq);
    \draw[flow] (lo) -- node[right,font=\scriptsize] {$e^{j\phi_n}$} (mix);

    \draw[flow,dashed] (est.north) -- ++(0,5mm)
      -| node[pos=0.25,above,font=\scriptsize] {$\hat{\phi}_n$ feedback} (fft.north);

    \begin{scope}[on background layer]
      \node[draw,dashed,rounded corners=2pt,fit=(arr)(ps),inner sep=2.5mm,
            label={[font=\scriptsize]below:analog domain}] {};
      \node[draw,dashed,rounded corners=2pt,fit=(fft)(est)(eq),inner sep=2.5mm,
            label={[font=\scriptsize]below:digital baseband}] {};
    \end{scope}
  \end{tikzpicture}
  \caption{Receiver signal chain. A single local oscillator drives the whole
  analog combining network, so its phase noise multiplies the beamformed signal
  and is common to every spatial stream. The estimator observes the
  frequency-domain samples of Eq.~\eqref{eq:freq-domain} and returns both a
  channel estimate and a phase trajectory.}
  \label{fig:chain}
\end{figure*}
